% ─────────────────────────────────────────────────────────────────────────────
This appendix provides the key code listings for the BookForMe system.  Full source
code is available at \url{https://github.com/JazibWaqas/JHAT}.

% ─────────────────────────────────────────────────────────────────────────────
\section{LangGraph Agent State Definition}
\label{app:code-langgraph}

The following listing shows the core LangGraph state schema and node graph
construction for the AI Receptionist.

\begin{lstlisting}[style=pythonstyle, caption={LangGraph agent state schema and graph construction (agent/graph.py).}]
from langgraph.graph import StateGraph, END
from typing import TypedDict, Optional, List
from enum import Enum

class BookingIntent(str, Enum):
    INQUIRY              = "inquiry"
    INFO                 = "info"
    TRANSACTION_CONFIRM  = "transaction_confirm"
    GREETING             = "greeting"
    UNKNOWN              = "unknown"

class AgentState(TypedDict):
    # Conversation context
    session_id:       str
    vendor_id:        str
    user_phone:       str
    messages:         List[dict]          # full conversation history

    # Extracted slot entities
    intent:           Optional[BookingIntent]
    service:          Optional[str]
    booking_date:     Optional[str]       # ISO 8601 date string
    booking_time:     Optional[str]       # HH:MM 24-hour
    duration_minutes: Optional[int]
    budget_max:       Optional[float]
    location:         Optional[str]

    # Booking state
    available_slots:  Optional[List[dict]]
    selected_slot:    Optional[dict]
    booking_id:       Optional[str]
    hold_expires_at:  Optional[str]

    # Payment state
    payment_required: Optional[float]
    screenshot_url:   Optional[str]
    ocr_amount:       Optional[float]

    # Control flow
    clarification_needed: bool
    error_message:    Optional[str]

def build_agent_graph() -> StateGraph:
    """Build and compile the cyclic LangGraph state machine."""
    graph = StateGraph(AgentState)

    # Register nodes
    graph.add_node("guardrails",        guardrails_node)
    graph.add_node("classify_intent",   classify_intent_node)
    graph.add_node("extract_slots",     extract_slots_node)
    graph.add_node("validate_state",    validate_state_node)
    graph.add_node("query_availability",query_availability_node)
    graph.add_node("execute_booking",   execute_booking_node)
    graph.add_node("request_payment",   request_payment_node)
    graph.add_node("verify_ocr",        verify_ocr_node)
    graph.add_node("notify_vendor",     notify_vendor_node)

    # Entry point
    graph.set_entry_point("guardrails")

    # Linear edges
    graph.add_edge("guardrails",          "classify_intent")
    graph.add_edge("classify_intent",     "extract_slots")
    graph.add_edge("extract_slots",       "validate_state")

    # Conditional: slots complete or need clarification
    graph.add_conditional_edges("validate_state", route_after_validation, {
        "clarify":    "classify_intent",   # loop back for clarification
        "available":  "query_availability"
    })

    # Conditional: slot available or not
    graph.add_conditional_edges("query_availability", route_after_availability, {
        "unavailable": END,                # send alternatives, end turn
        "available":   "execute_booking"
    })

    # Booking and payment flow
    graph.add_edge("execute_booking",     "request_payment")

    # Conditional: OCR pass or fail (loops until valid screenshot)
    graph.add_conditional_edges("verify_ocr", route_after_ocr, {
        "invalid":  "request_payment",    # loop: ask for new screenshot
        "valid":    "notify_vendor"
    })

    graph.add_edge("notify_vendor", END)

    return graph.compile()
\end{lstlisting}

% ─────────────────────────────────────────────────────────────────────────────
\section{Optimistic Concurrency Control Implementation}
\label{app:code-occ}

\begin{lstlisting}[style=pythonstyle, caption={OCC slot locking via Firestore transaction (agent/booking.py).}]
import firebase_admin
from firebase_admin import firestore
from datetime import datetime, timedelta, timezone
import uuid

db = firestore.client()
HOLD_DURATION_MINUTES = 10

def create_booking_hold(
    service_id: str,
    slot_start: str,
    user_phone: str,
    vendor_id:  str
) -> dict:
    """
    Atomically lock a time slot using Optimistic Concurrency Control.
    Returns booking record on success, raises ConflictError on failure.
    """
    slot_ref = (
        db.collection("vendors")
          .document(vendor_id)
          .collection("slots")
          .document(f"{service_id}_{slot_start}")
    )

    booking_id = f"BFM-{uuid.uuid4().hex[:8].upper()}"
    hold_expiry = datetime.now(timezone.utc) + timedelta(
                      minutes=HOLD_DURATION_MINUTES)

    @firestore.transactional
    def run_transaction(transaction):
        snapshot = slot_ref.get(transaction=transaction)

        if not snapshot.exists:
            raise ValueError(f"Slot {slot_start} does not exist.")

        data = snapshot.to_dict()
        status = data.get("status", "AVAILABLE")

        # Check for active HOLD that has not expired
        if status == "HOLD":
            expiry = data.get("hold_expires_at")
            if expiry and expiry > datetime.now(timezone.utc):
                raise ConflictError(
                    f"Slot {slot_start} is currently held by another user.")

        if status == "CONFIRMED":
            raise ConflictError(
                f"Slot {slot_start} is already booked.")

        # Atomic write: transition to HOLD
        transaction.update(slot_ref, {
            "status":          "HOLD",
            "booking_id":      booking_id,
            "user_phone":      user_phone,
            "hold_expires_at": hold_expiry,
        })

    transaction = db.transaction()
    run_transaction(transaction)

    return {
        "booking_id":      booking_id,
        "status":          "HOLD",
        "hold_expires_at": hold_expiry.isoformat(),
        "service_id":      service_id,
        "slot_start":      slot_start,
    }


class ConflictError(Exception):
    """Raised when a slot is concurrently locked by another user."""
    pass
\end{lstlisting}

% ─────────────────────────────────────────────────────────────────────────────
\section{Intent Classification (DistilBERT Inference)}
\label{app:code-nlu}

\begin{lstlisting}[style=pythonstyle, caption={Intent classification using fine-tuned DistilBERT (agent/nlu.py).}]
from transformers import pipeline
from agent.graph import BookingIntent
import os

# Load fine-tuned DistilBERT model from local path or HuggingFace Hub
_NLU_MODEL_PATH = os.getenv("NLU_MODEL_PATH", "./models/distilbert-booking")
CONFIDENCE_THRESHOLD = 0.65   # below this, route to Gemini fallback

_classifier = None

def get_classifier():
    global _classifier
    if _classifier is None:
        _classifier = pipeline(
            "text-classification",
            model=_NLU_MODEL_PATH,
            top_k=None,                  # return scores for all classes
            truncation=True,
            max_length=128,
        )
    return _classifier


def classify_intent(text: str) -> tuple[BookingIntent, float]:
    """
    Classify a user message into one of the five BookingIntent classes.

    Returns:
        (intent, confidence) where intent is a BookingIntent enum value
        and confidence is the softmax probability of the top class.

    If top confidence < CONFIDENCE_THRESHOLD, returns
    (BookingIntent.UNKNOWN, confidence) to trigger LLM fallback.
    """
    clf = get_classifier()
    results = clf(text)[0]                # list of {label, score}

    # Sort by score descending
    results_sorted = sorted(results, key=lambda x: x["score"], reverse=True)
    top = results_sorted[0]

    label_map = {
        "LABEL_0": BookingIntent.GREETING,
        "LABEL_1": BookingIntent.INQUIRY,
        "LABEL_2": BookingIntent.INFO,
        "LABEL_3": BookingIntent.TRANSACTION_CONFIRM,
        "LABEL_4": BookingIntent.UNKNOWN,
    }

    intent = label_map.get(top["label"], BookingIntent.UNKNOWN)
    confidence = top["score"]

    if confidence < CONFIDENCE_THRESHOLD:
        return BookingIntent.UNKNOWN, confidence

    return intent, confidence
\end{lstlisting}

% ─────────────────────────────────────────────────────────────────────────────
\section{OCR Payment Verification}
\label{app:code-ocr}

\begin{lstlisting}[style=pythonstyle, caption={OCR payment amount extraction (services/ocr\_service.py).}]
import re
import requests
from PIL import Image
import pytesseract
import io

PKR_PATTERNS = [
    r'Rs\.?\s*([\d,]+)',
    r'PKR\s*([\d,]+)',
    r'Amount[:\s]+([\d,]+)',
    r'([\d,]+)\s*(?:rupees?|PKR|Rs)',
]

def extract_payment_amount(image_url: str) -> float | None:
    """
    Download payment screenshot from Firebase Storage and extract
    the transaction amount using Tesseract OCR.

    Returns:
        Extracted amount as float, or None if extraction fails.
    """
    # Download image from signed URL
    response = requests.get(image_url, timeout=10)
    response.raise_for_status()

    # Open with Pillow and run Tesseract
    image = Image.open(io.BytesIO(response.content))
    raw_text = pytesseract.image_to_string(image, config="--psm 6")

    # Try each PKR pattern in order of specificity
    for pattern in PKR_PATTERNS:
        match = re.search(pattern, raw_text, re.IGNORECASE)
        if match:
            amount_str = match.group(1).replace(",", "")
            try:
                return float(amount_str)
            except ValueError:
                continue

    return None   # OCR could not extract a valid amount


def verify_payment(
    image_url:       str,
    required_amount: float,
    tolerance:       float = 0.01   # 1% tolerance for rounding
) -> tuple[bool, float | None]:
    """
    Verify that the payment screenshot shows the required advance amount.

    Returns:
        (is_valid, extracted_amount)
    """
    extracted = extract_payment_amount(image_url)

    if extracted is None:
        return False, None

    is_valid = abs(extracted - required_amount) <= (required_amount * tolerance)
    return is_valid, extracted
\end{lstlisting}

% ─────────────────────────────────────────────────────────────────────────────
\section{Elo Rating Update}
\label{app:code-elo}

\begin{lstlisting}[style=pythonstyle, caption={Elo rating calculation and update (social/elo.py).}]
K_FACTOR = 32   # development coefficient for all players

def expected_score(rating_a: float, rating_b: float) -> float:
    """Compute expected score for player A against player B."""
    return 1.0 / (1.0 + 10 ** ((rating_b - rating_a) / 400.0))


def update_elo(
    rating_a: float,
    rating_b: float,
    score_a:  float   # 1.0 = win, 0.5 = draw, 0.0 = loss
) -> tuple[float, float]:
    """
    Compute new Elo ratings for both players after a match.

    Returns:
        (new_rating_a, new_rating_b)
    """
    ea = expected_score(rating_a, rating_b)
    eb = expected_score(rating_b, rating_a)
    score_b = 1.0 - score_a

    new_a = rating_a + K_FACTOR * (score_a - ea)
    new_b = rating_b + K_FACTOR * (score_b - eb)

    return round(new_a, 1), round(new_b, 1)


def matchmaking_window(wait_seconds: int, base: int = 200,
                       expansion: int = 50) -> int:
    """
    Compute the Elo search window that expands with wait time.
    Expands by `expansion` points for every 60 seconds waiting.
    """
    return base + expansion * (wait_seconds // 60)
\end{lstlisting}

% ─────────────────────────────────────────────────────────────────────────────
\section{WhatsApp Webhook Handler}
\label{app:code-webhook}

\begin{lstlisting}[style=pythonstyle, caption={FastAPI WhatsApp webhook endpoint (api/whatsapp\_webhook.py).}]
from fastapi import APIRouter, Request, HTTPException, Query
from agent.graph   import build_agent_graph
from agent.session import load_session, save_session
import os, hashlib, hmac, json

router = APIRouter()
_agent = build_agent_graph()

VERIFY_TOKEN  = os.environ["WHATSAPP_VERIFY_TOKEN"]
APP_SECRET    = os.environ["WHATSAPP_APP_SECRET"]

# -- Webhook verification (GET) -----------------------------------------------
@router.get("/webhook")
async def verify_webhook(
    hub_mode:        str = Query(alias="hub.mode"),
    hub_challenge:   str = Query(alias="hub.challenge"),
    hub_verify_token:str = Query(alias="hub.verify_token"),
):
    if hub_mode == "subscribe" and hub_verify_token == VERIFY_TOKEN:
        return int(hub_challenge)
    raise HTTPException(status_code=403, detail="Forbidden")

# -- Inbound message handler (POST) -------------------------------------------
@router.post("/webhook")
async def receive_message(request: Request):
    # Verify Meta signature
    body_bytes = await request.body()
    signature  = request.headers.get("X-Hub-Signature-256", "")
    expected   = "sha256=" + hmac.new(
        APP_SECRET.encode(), body_bytes, hashlib.sha256).hexdigest()
    if not hmac.compare_digest(signature, expected):
        raise HTTPException(status_code=403, detail="Invalid signature")

    payload = json.loads(body_bytes)

    # Extract message details from Meta webhook payload
    for entry in payload.get("entry", []):
        for change in entry.get("changes", []):
            value = change.get("value", {})
            for msg in value.get("messages", []):
                user_phone = msg.get("from")
                vendor_id  = value.get("metadata", {}).get("phone_number_id")

                # Load or initialise conversation session
                state = await load_session(user_phone, vendor_id)

                # Append new user message
                if msg.get("type") == "text":
                    state["messages"].append({
                        "role": "user",
                        "content": msg["text"]["body"]
                    })
                elif msg.get("type") == "image":
                    state["screenshot_url"] = msg["image"]["id"]
                    state["messages"].append({
                        "role": "user",
                        "content": "[Payment screenshot uploaded]"
                    })

                # Run agent graph
                updated_state = await _agent.ainvoke(state)

                # Persist updated session
                await save_session(user_phone, vendor_id, updated_state)

    return {"status": "ok"}
\end{lstlisting}

\bigskip
\noindent The complete source code, including the React Native frontend, vendor
dashboard, training scripts, and deployment configuration, is available at:

\begin{center}
  \url{https://github.com/JazibWaqas/JHAT}
\end{center}

%%% Local Variables:
%%% mode: latex
%%% TeX-master: "../report"
%%% End:
